\documentclass[11pt,a4paper]{article}
\usepackage[margin=21mm,headheight=15pt]{geometry}
\usepackage[T1]{fontenc}
\usepackage{lmodern,amsmath,amssymb,booktabs,tabularx,enumitem,xcolor,fancyhdr}
\usepackage[hidelinks]{hyperref}
\definecolor{accent}{RGB}{0,114,178}
\setlength{\parindent}{0pt}
\setlength{\parskip}{6pt}
\setlist{nosep,leftmargin=*}
\pagestyle{fancy}
\fancyhf{}
\fancyhead[L]{\small Econometrics 25117 | Instructor guide}
\fancyhead[R]{\small 2026--2027}
\fancyfoot[L]{\small Private teaching notes}
\fancyfoot[R]{\thepage}
\newcommand{\heading}[1]{\section*{\color{accent}#1}}
\newcommand{\cue}[2]{\par\noindent\textbf{#1}\quad #2\par}
\setlength{\emergencystretch}{1em}
\newcommand{\E}{\mathbb{E}}
\newcommand{\Var}{\operatorname{Var}}
\newcommand{\Cov}{\operatorname{Cov}}
\begin{document}
\begin{center}{\Large\bfseries Simple regression and OLS}\end{center}
\textbf{Slide route:} Lecture3v2.tex: regression line through large-sample distribution. Stock and Watson Chapter 4.
\heading{Session 4 | A line, an estimator, an assumption}
\textbf{Outcomes:} Compute an OLS line, distinguish errors from residuals, interpret slope units, and state the role of conditional mean zero.
\cue{Opening}{A line can summarize an association. What extra argument would let us call its slope a causal effect? Keep that question on the board throughout the algebra.}
\heading{100-minute route}
\begin{tabularx}{\textwidth}{@{}rX@{}}\toprule
0--10 & Retrieve covariance and conditional expectation.\\
10--25 & Population model, fitted line, errors versus residuals.\\
25--45 & OLS objective and the three-observation example.\\
45--55 & Interpretation and pause.\\
55--70 & $R^2$, SER, fit versus causality.\\
70--90 & Causal assumptions; unbiasedness and sampling variation.\\
90--100 & Exit ticket; hand off inference to Lecture 4.\\\bottomrule
\end{tabularx}
\heading{Board example | Compute, do not memorize}
\cue{Use with slides}{Lecture3v2 slides 4--9: The regression line, Population Linear Regression Model, and OLS Estimator. Begin the three-point calculation when slide 7 introduces OLS; use slides 8--9 for the first-order conditions and residual identities.}
Use synthetic observations $(x,y)=(1,2),(2,3),(3,5)$.
\[
\bar x=2,\quad \bar y=10/3,\quad
\hat\beta_1=\frac{\sum(x_i-\bar x)(y_i-\bar y)}
{\sum(x_i-\bar x)^2}=\frac3{2}=1.5.
\]
\[
\hat\beta_0=\bar y-\hat\beta_1\bar x=1/3.
\]
Predictions are $11/6,10/3,29/6$; residuals are $1/6,-1/3,1/6$. Their sum is zero, and $\sum x_i\hat u_i=0$.
\cue{Derivation cue}{Differentiate $\sum(y_i-b_0-b_1x_i)^2$ with respect to both parameters. The first equation puts the fitted line through $(\bar x,\bar y)$; substitute it into the second to obtain the slope formula.}
\cue{Distinction}{Population $u_i$ is unobserved. Sample $\hat u_i=y_i-\hat y_i$ is computed. OLS residual orthogonality is an algebraic identity, not evidence of exogeneity.}

\newpage
\heading{Use the same example | Fit and interpretation}
\cue{Use with slides}{Lecture3v2 slides 11--15: Interpretation, Goodness of Fit, $R^2$, SER, and Goodness of Fit again. Use the same three observations to compute $R^2$ on slide 13 and SER on slide 14.}
\[
SSR=\sum\hat u_i^2=1/6,\qquad
TSS=\sum(y_i-\bar y)^2=14/3.
\]
Here SSR means residual sum of squares; always define the notation because some texts use it for regression sum of squares.
\[
R^2=1-\frac{1/6}{14/3}=\frac{27}{28}\simeq.9643,
\qquad SER=\sqrt{\frac{1/6}{3-2}}=\sqrt{1/6}\simeq.408.
\]
These three observations demonstrate arithmetic, not reliable asymptotic inference.
\cue{Units exercise}{If $y$ is a test score and $x$ is class size, slope 1.5 means an associated 1.5-point difference per additional student. If $x$ is divided by 10, the numerical slope multiplies by 10. Predictions do not change.}
\heading{Causality and sampling}
\cue{Use with slides}{Lecture3v2 slides 16--26: conditions for causality, conditional mean zero, i.i.d. draws, outliers, and OLS sample properties. Use the bias expression beside slides 18--19 before moving to large-sample distribution.}
For $Y_i=\beta_0+\beta_1X_i+u_i$, a zero conditional mean $\E(u_i\mid X_i)=0$ says omitted determinants average to zero at each regressor value. Richer districts may have both different class sizes and different resources, violating this condition.
\[
\hat\beta_1-\beta_1=
\frac{\sum_i(X_i-\bar X)u_i}{\sum_i(X_i-\bar X)^2}.
\]
Use the expression to explain why variation in $X$ is needed and why correlation between $X$ and omitted determinants is a problem. Discuss random sampling and finite-moment conditions before invoking large-sample normality.
\cue{Side question}{Does a high $R^2$ establish causality? No. Does a low $R^2$ rule out a causal effect? No. Fit and identification answer different questions.}
\cue{Exit answers}{Residuals sum to zero when an intercept is included. That fact cannot verify $\E(u\mid X)=0$. The intercept may have no sensible substantive interpretation when $X=0$ is outside the observed range.}
\cue{If short / preparation}{Assign the detailed asymptotic variance algebra for review. Preserve the numerical OLS calculation and exogeneity discussion. Read Chapter 5 before session 5.}
\textbf{Source:} Lecture3v2.tex; synthetic three-observation calculation.

\end{document}
